\documentclass[10pt,a4paper]{amsart}
\usepackage{amsmath,amssymb,amsthm,graphicx}
% For commutative diagrams we use Xy-pic
\usepackage[all]{xy}

% We use 2cell for 2-commutative diagrams.
\xyoption{2cell}
% Theorem environments.
%
\theoremstyle{plain}
\newtheorem{theorem}[subsection]{Theorem}
\newtheorem{proposition}[subsection]{Proposition}
\newtheorem{lemma}[subsection]{Lemma}

\theoremstyle{definition}
\newtheorem{definition}[subsection]{Definition}
\newtheorem{example}[subsection]{Example}
\newtheorem{exercise}[subsection]{Exercise}
\newtheorem{situation}[subsection]{Situation}

\theoremstyle{remark}
\newtheorem{remark}[subsection]{Remark}
\newtheorem{remarks}[subsection]{Remarks}

\numberwithin{equation}{subsection}

% Macros
%
\def\lim{\mathop{\mathrm{lim}}\nolimits}
\def\colim{\mathop{\mathrm{colim}}\nolimits}
\def\Spec{\mathop{\mathrm{Spec}}}
\def\Hom{\mathop{\mathrm{Hom}}\nolimits}
\def\Ext{\mathop{\mathrm{Ext}}\nolimits}
\def\SheafHom{\mathop{\mathcal{H}\!\mathit{om}}\nolimits}
\def\SheafExt{\mathop{\mathcal{E}\!\mathit{xt}}\nolimits}
\def\Sch{\mathit{Sch}}
\def\Mor{\mathop{\mathrm{Mor}}\nolimits}
\def\Ob{\mathop{\mathrm{Ob}}\nolimits}
\def\Sh{\mathop{\mathit{Sh}}\nolimits}
\def\NL{\mathop{N\!L}\nolimits}
\def\CH{\mathop{\mathrm{CH}}\nolimits}
\def\proetale{{pro\text{-}\acute{e}tale}}
\def\etale{{\acute{e}tale}}
\def\QCoh{\mathit{QCoh}}
\def\Ker{\mathop{\mathrm{Ker}}}
\def\Im{\mathop{\mathrm{Im}}}
\def\Coker{\mathop{\mathrm{Coker}}}
\def\Coim{\mathop{\mathrm{Coim}}}

% Boxtimes
%
\DeclareMathSymbol{\boxtimes}{\mathbin}{AMSa}{"02}

%
% Macros for moduli stacks/spaces
%
\def\QCohstack{\mathcal{QC}\!\mathit{oh}}
\def\Cohstack{\mathcal{C}\!\mathit{oh}}
\def\Spacesstack{\mathcal{S}\!\mathit{paces}}
\def\Quotfunctor{\mathrm{Quot}}
\def\Hilbfunctor{\mathrm{Hilb}}
\def\Curvesstack{\mathcal{C}\!\mathit{urves}}
\def\Polarizedstack{\mathcal{P}\!\mathit{olarized}}
\def\Complexesstack{\mathcal{C}\!\mathit{omplexes}}
% \Pic is the operator that assigns to X its picard group, usage \Pic(X)
% \Picardstack_{X/B} denotes the Picard stack of X over B
% \Picardfunctor_{X/B} denotes the Picard functor of X over B
\def\Pic{\mathop{\mathrm{Pic}}\nolimits}
\def\Picardstack{\mathcal{P}\!\mathit{ic}}
\def\Picardfunctor{\mathrm{Pic}}
\def\Deformationcategory{\mathcal{D}\!\mathit{ef}}
\newenvironment{reference}{}{}
\newenvironment{slogan}{}{}
\newenvironment{history}{}{}
\usepackage{cleveref}
\crefname{originallemma}{original lemma}{original lemmas}
\Crefname{originallemma}{Original lemma}{Original lemmas}
\makeatletter
\newlabel{lemma-divisor-special-fiber}{{0C5Z}{source}{}{sourcepointer-1}{}}
\newlabel{lemma-divisor-special-fiber@cref}{{[originallemma][0][]0C5Z}{[1][0][]source}{}{}{}}
\makeatother
\crefname{originallemma}{original lemma}{original lemmas}
\Crefname{originallemma}{Original lemma}{Original lemmas}
\makeatletter
\newlabel{lemma-regular-model-pic}{{0C63}{source}{}{sourcepointer-2}{}}
\newlabel{lemma-regular-model-pic@cref}{{[originallemma][0][]0C63}{[1][0][]source}{}{}{}}
\makeatother
\crefname{originallemma}{original lemma}{original lemmas}
\Crefname{originallemma}{Original lemma}{Original lemmas}
\makeatletter
\newlabel{lemma-properties-form}{{0C66}{source}{}{sourcepointer-3}{}}
\newlabel{lemma-properties-form@cref}{{[originallemma][0][]0C66}{[1][0][]source}{}{}{}}
\makeatother
\crefname{originallemma}{original lemma}{original lemmas}
\Crefname{originallemma}{Original lemma}{Original lemmas}
\makeatletter
\newlabel{lemma-multiple-fibre-normal-bundle}{{0C67}{source}{}{sourcepointer-4}{}}
\newlabel{lemma-multiple-fibre-normal-bundle@cref}{{[originallemma][0][]0C67}{[1][0][]source}{}{}{}}
\makeatother
\title{An invertible sheaf trivial on the curve and of degree zero on every component becomes trivial after raising to the gcd of the multiplicities}
\author{The Stacks Project Authors}
\date{stacks-project:models.tex:local-snapshot-20261012; GNU FDL 1.2 or later}
\begin{document}
\maketitle
\noindent\textit{Editorial note.} The counted unit is the Stacks Project Lemma lemma-characterize-trivial at source lines 5886-5895 together with its proof at source lines 5897-5909, a straight adjacent pairing in the chapter file. The chapter is the upstream file models.tex, which is the file name recorded in the source record; the pool directory carries the longer name stacks-models and is not used as the chapter name anywhere in this card.\par\medskip
\noindent\textit{Editorial note.} The retained statement opens with the words In Situation followed by a reference to situation-regular-model, so the situation environment that defines the setting is retained as a context span at source lines 4628-4637. It is context, not a counted unit, and it does not overlap the counted interval.\par\medskip
\noindent\textit{Editorial note.} This chapter's candidate index records the statement body and the proof body without their closing markers, so the retained spans used here are the full environments, exactly as in the earlier cards of this chapter.\par\medskip
\noindent\textit{Editorial note.} The retained proof names three results of the same chapter, for the exact sequence of the regular model, for the proportionality of a combination with vanishing self-intersection to the multiplicity vector, and for the torsion order of the divided special fibre. All three are declared as non-load-bearing navigation pointers with their stable Stacks tags looked up in the shared tag index; the retained situation additionally names one lemma of this chapter, also declared as a pointer.\par\medskip
\noindent\textit{Editorial note.} No citation occurs in the retained spans, so no bibliography entry is supplied, and the retained preamble is the chapter preamble's Xy-pic setup at lines 14-18, the theorem declarations and mathematical macros at lines 154-217, and the three no-op wrapper environments.\par\medskip
\noindent\textit{Allowed standard prerequisites (editorial).}\begin{itemize}
\item The exact sequence of the regular model, the positivity of the intersection form and the divided special fibre. For a regular model an invertible sheaf that is trivial on the generic fibre is attached to a divisor supported on the special fibre and determined by it; the intersection form on the components is negative definite on the subspace of combinations orthogonal to the multiplicity vector, so a combination of components whose self-intersection number vanishes is a rational multiple of the multiplicity vector; and the divisor obtained by dividing the special fibre by the gcd of the multiplicities is torsion of order dividing that gcd. The author's own lemmas of this section are the allowed inputs.. The counted lemma is the characterisation used to control the kernel in the h-torsion sequence of this section and to make the divisor of the divided special fibre explicit.
\end{itemize}\medskip
\ifdefined\hypertarget\hypertarget{sourcepointer-1}{}\fi
\noindent\textit{Source reference.} Original lemma 0C5Z: focused-seeds/AGP/sources/stacks-snapshot-20261011/tags.txt, source lines 15279--15279. Navigation pointer inside the retained situation, which names it when it writes the special fibre as a sum of its components with multiplicities. The number is the stable Stacks tag of the target, looked up in the shared tag index.\par\medskip
\ifdefined\hypertarget\hypertarget{sourcepointer-2}{}\fi
\noindent\textit{Source reference.} Original lemma 0C63: focused-seeds/AGP/sources/stacks-snapshot-20261011/tags.txt, source lines 15283--15283. Pointer to the exact sequence of the regular model, used by the retained proof to write the sheaf as an invertible sheaf attached to a combination of components; that exact sequence is counted by a separate card of this lane and is not reproduced here. The number is the stable Stacks tag of the target, looked up in the shared tag index.\par\medskip
\ifdefined\hypertarget\hypertarget{sourcepointer-3}{}\fi
\noindent\textit{Source reference.} Original lemma 0C66: focused-seeds/AGP/sources/stacks-snapshot-20261011/tags.txt, source lines 15286--15286. Pointer to the author's statement that a combination of components whose self-intersection number vanishes is a rational multiple of the multiplicity vector; the cited statement is an allowed input of the same section and is not reproduced here. The number is the stable Stacks tag of the target, looked up in the shared tag index.\par\medskip
\ifdefined\hypertarget\hypertarget{sourcepointer-4}{}\fi
\noindent\textit{Source reference.} Original lemma 0C67: focused-seeds/AGP/sources/stacks-snapshot-20261011/tags.txt, source lines 15287--15287. Pointer to the torsion statement for the divided special fibre, used by the retained proof to conclude that the d-th tensor power is trivial; that statement is counted by a separate card of this lane and is not reproduced here. The number is the stable Stacks tag of the target, looked up in the shared tag index.\par\medskip
\begin{situation}
\label{situation-regular-model}
Let $R$ be a discrete valuation ring with fraction field $K$,
residue field $k$, and uniformizer $\pi$.
Let $C$ be a smooth projective curve over $K$ with $H^0(C, \mathcal{O}_C) = K$.
Let $X$ be a regular proper model of $C$.
Let $C_1, \ldots, C_n$ be the irreducible components of the special
fibre $X_k$. Write $X_k = \sum m_i C_i$ as in
Lemma \ref{lemma-divisor-special-fiber}.
\end{situation}

\begin{lemma}
\label{lemma-characterize-trivial}
In Situation \ref{situation-regular-model} let $d = \gcd(m_1, \ldots, m_n)$.
If $\mathcal{L}$ is an invertible $\mathcal{O}_X$-module which
\begin{enumerate}
\item restricts to the trivial invertible module on $C$, and
\item has degree $0$ on each $C_i$,
\end{enumerate}
then $\mathcal{L}^{\otimes d} \cong \mathcal{O}_X$.
\end{lemma}

\begin{proof}
By Lemma \ref{lemma-regular-model-pic} we have
$\mathcal{L} \cong \mathcal{O}_X(\sum a_i C_i)$ for some
$a_i \in \mathbf{Z}$. The degree of $\mathcal{L}|_{C_j}$
is $\sum_j a_j(C_i \cdot C_j)$. In particular
$(\sum a_i C_i \cdot \sum a_i C_i) = 0$.
Hence we see from Lemma \ref{lemma-properties-form}
that $(a_1, \ldots, a_n) = q(m_1, \ldots, m_n)$ for some
$q \in \mathbf{Q}$. Thus $\mathcal{L} = \mathcal{O}_X(lD)$
for some $l \in \mathbf{Z}$ where $D = \sum (m_i/d) C_i$ is as in
Lemma \ref{lemma-multiple-fibre-normal-bundle}
and we conclude.
\end{proof}

\end{document}
